\scanpage{555}

thus $\{1,\tau'\}$ acts on these groups, and we can consider
\[
\tag{34}
\left\{\begin{aligned}
N_\rho&=\{1,\tau'\}\mathbin{.}Z_\rho,\qquad\deleted{\unsure{N}}\\
N_\sigma&=\{1,\tau'\}\mathbin{.}Z_\sigma
\end{aligned}\right.
\qquad(\text{products }\tfrac12\text{ direct}).
\]

We have canonical exact sequences
\[
\tag{35}
\diagram{\begin{tikzcd}
1\arrow[r]&\unsure{Z_\rho|L_\rho}\arrow[r]&N_\rho\arrow[r,"\varepsilon_\rho"]&\{\pm1\}\arrow[r]&1\\
1\arrow[r]&Z_\sigma\arrow[r]&N_\sigma\arrow[r,"\varepsilon_\sigma"]&\{\pm1\}\arrow[r]&1
\end{tikzcd}},
\]
and a commutative diagram of canonical homomorphisms
\[
\tag{36}
\diagram{\begin{tikzcd}[row sep=small,column sep=large]
&\mathbb D_\rho\arrow[r]&N_\rho\arrow[dr,"i_\rho"]&\\
\{1,\tau'\}\times\{\pm1\}\arrow[ur,hook]\arrow[dr,hook]&&&G_{\rho,\sigma}\\
&\mathbb D_\sigma\arrow[r]&N_\sigma\arrow[ur,"\unsure{i_\sigma}"']&
\end{tikzcd}}
\]
where
\[
\tag{37}
\left\{\begin{aligned}
\mathbb D_\rho&\simeq\mathbb D_6\subset\operatorname{Gl}(2,\mathbb Z)
\quad\text{generated by }\rho,\tau'\\
\mathbb D_\sigma&\simeq\mathbb D_4\subset\operatorname{Gl}(2,\mathbb Z)
\quad\text{\rule{2.7cm}{0.35pt}}\ \sigma,\tau'\\
\{1,\tau'\}\times\{\pm1\}&=\mathbb D_\rho\cap\mathbb D_\sigma
\quad\text{in }\operatorname{Gl}(2,\mathbb Z),
\end{aligned}\right.
\]
\deleted{\unsure{hence} \unsure{the homomorphisms} $\mathbb D_\rho\longrightarrow N_\rho$}
\[
\tag{38}
\left\{\begin{aligned}
\mathbb D_\rho&=\{1,\tau'\}\mathbin{.}L_\rho
\quad(\text{at least if }\rho\text{ of order }6\text{ in }\mathcal G)\\
\mathbb D_\sigma&=\{1,\tau'\}\mathbin{.}L_\sigma
\quad(\text{\rule{1.5cm}{0.35pt}}\ \sigma\text{ of order }4\ \text{\rule{1.4cm}{0.35pt}}),
\end{aligned}\right.
\]
and the homomorphisms
\[
\diagram{\begin{tikzcd}
\mathbb D_\rho\arrow[r]&N_\rho,
\end{tikzcd}}
\qquad
\diagram{\begin{tikzcd}
\mathbb D_\sigma\arrow[r]&N_\sigma
\end{tikzcd}}
\]
are obvious. The homomorphisms $i_\rho,i_\sigma$ too, of course.

Let us recall that
\[
\tag{39}
\left\{\begin{aligned}
\deleted{\unsure{i'_\rho}}i_\rho\text{ injective}
&\Longleftrightarrow\tau'\notin Z_\rho
\Longleftrightarrow\rho^2\ne1
\Longleftrightarrow\rho\ne-1
\Longleftrightarrow\mathcal G^+\text{ is }\unsure{\deleted{\text{isomorphic}}|\text{isomorphic}}
\text{ to a quotient of }\mathbb Z/4\mathbb Z.\\
i_\sigma\text{ injective}
&\Longleftrightarrow\tau'\notin Z_\sigma
\Longleftrightarrow\sigma^2\ne1\quad
\deleted{\unsure{\Longleftrightarrow\sigma\ne\pm1}}
\end{aligned}\right.
\]
\mspage{554}

{\small
\begin{center}
554\\
\texttt{https://grothendieck.umontpellier.fr}\\
Université de Montpellier, UMR CNRS 5149, Case courrier 051, Place Eugène Bataillon\\
34095 Montpellier cedex 5 - France
\end{center}
}

\scanpage{556}
\mspage{581}

Letting $G_{\rho,\sigma}$ act on $\mathcal G^+$, a normal subgroup,
via inner automorphisms, we see that every element of $G_{\rho,\sigma}$ transforms
$L_\rho\subset\mathcal G^+$ into a \inserted{$\mathcal G^+$-}conjugate subgroup,
thus is \unsure{congruent} modulo \deleted{$\unsure{\operatorname{Norm}_{\mathcal G^+}(L_\rho)}$}
$\mathcal G^+$ to an element of
$\operatorname{Norm}_{\unsure{G|\mathcal G}}(L_\rho)=\unsure{\operatorname{Im}}i_\rho(N_\rho)$,
and likewise for $L_\sigma$. In other words, \inserted{setting}
\inserted{\[
\tag{40}
\begin{aligned}
N_\rho^*&\overset{\mathrm{dfn}}=\{x\in N_\rho\mid\varepsilon_\rho(x)=\varepsilon'_G(i_\rho(x))\},\qquad
N_\sigma^*\overset{\mathrm{dfn}}=\{x\in N_\sigma\mid\varepsilon_\sigma(x)=\varepsilon_G(i_\sigma(x))\}
\end{aligned}
\]}
the composite homomorphisms
\[
\tag{\unsure{40}}
\left\{\begin{array}{ll}
\diagram{\begin{tikzcd}[column sep=large]
N_\rho^*\arrow[r,"i_\rho"]&G_{\rho,\sigma}\arrow[r,"\delta_G"]&\gamma_{\rho,\sigma}
\end{tikzcd}},&\delta_\rho\overset{\mathrm{dfn}}=\delta_G\circ i_\rho\\[1.5em]
\diagram{\begin{tikzcd}[column sep=large]
N_\sigma^*\arrow[r,"i_\sigma"]&G_{\rho,\sigma}\arrow[r,"\delta_G"]&\gamma_{\rho,\sigma}
\end{tikzcd}},&\delta_\sigma\overset{\mathrm{dfn}}=\delta_G\circ i_\sigma
\end{array}\right.
\]
epimorphisms.
\deleted{We \unsure{see} \illegible{3}}
\inserted{We thus have \deleted{\illegible{1}}}

\margin{If $g\in G_{\rho,\sigma}$, and
$\varepsilon=\varepsilon_G(g)$, $\varepsilon'=\varepsilon'_G(g)$, then
\[
\begin{aligned}
\operatorname{int}(g)(\rho)&\underset{\mathcal G^+}{\sim}\rho^{\varepsilon'}
=\unsure{\tau^{\nu'}(\rho)},\\
\operatorname{int}(g)(\sigma)&\underset{\mathcal G^+}{\sim}\sigma^\varepsilon
=\unsure{\tau^\nu(\rho)|\tau^\nu(\sigma)}.
\end{aligned}
\]
\deleted{\unsure{thus}}
(where $(-1)^\nu=\varepsilon$, $(-1)^{\nu'}=\varepsilon'$), \inserted{thus}
\deleted{$\exists g\in\unsure{\mathcal G^+}$ such that}
$\exists a=\tau^{\nu'}\mathbin{.}\unsure{a_0|ag}\in N_\rho$
(where $a_0=\unsure{\tau^\nu g}$) such that

a) $g\equiv\unsure{a|a_0}\ (\mathcal G^+)$ and

b) $\varepsilon_\rho(a)=\unsure{\varepsilon'|\varepsilon'_G}
=\unsure{\varepsilon'_G(g)|\varepsilon_G(g)}$ ---

(and idem for $N_\sigma$\unsure{\dots})

thus
\[
\begin{aligned}
\operatorname{Im}\bigl(N_\rho^{\deleted{\unsure{*}}(+1)}
\xrightarrow{\deleted{\unsure{\delta_\rho}}}\gamma_{\rho,\sigma}\bigr)
&\supset\gamma'_{\rho,\sigma}\deleted{\unsure{\varepsilon'}},\\
\operatorname{Im}\bigl(N_\rho^{(-1)}\longrightarrow\gamma_{\rho,\sigma}\bigr)
&\supset\gamma_{\rho,\sigma}\setminus\gamma'_{\rho,\sigma}.
\end{aligned}
\]
But this does \uline{not} prove that we have \unsure{commutativity}!
\unsure{hence probably}
}

\deleted{\unsure{They are not} \illegible{2} (cf. (39))}
\deleted{(\unsure{41}) a) If $\rho^2\ne1$, then we have commutativity in}
\[
\diagram{\deleted{\begin{tikzcd}[column sep=large,row sep=small]
N_\rho\arrow[rr,"\delta_\rho"]
\arrow[dr,"\varepsilon_\rho"']&&
\gamma_{\rho,\sigma}\arrow[dl,"\varepsilon'_\gamma"]\\
\unsure{\varepsilon_\rho}&\{\pm1\}&\\
&\unsure{\{\pm1\}}&\unsure{\{\pm1\}}
\end{tikzcd}}}
\]
\deleted{hence}
\inserted{commutative} diagram of \deleted{$\unsure{\{\pm1\}}$} exact sequences
\[
\tag{41}
\diagram{\begin{tikzcd}[column sep=large]
1\arrow[r]&Z_\rho\arrow[r]\arrow[d,"\text{epimorphism}"]&N_\rho^*\arrow[r,"\varepsilon_\rho"]\arrow[d,"\text{epimorphism}"]&\{\pm1\}\arrow[r]\arrow[d,equal]&1\\
1\arrow[r]&\gamma'_{\rho,\sigma}\arrow[r]&\unsure{\gamma_{\rho,\sigma}|Z_{\rho,\sigma}}\arrow[r,"\varepsilon'_\gamma"]&\{\pm1\}\arrow[r]&1
\end{tikzcd}}
\]
where we have set
\[
\tag{42}
\gamma'_{\rho,\sigma}=\operatorname{Ker}\bigl(
\diagram{\begin{tikzcd}
\gamma_{\rho,\sigma}\arrow[r,"\varepsilon'_\gamma"]&\{\pm1\}
\end{tikzcd}}\bigr);
\]

\deleted{b) If $\sigma^2\ne1$ (cf. (39)) then we have commutativity}
\inserted{Likewise}
\[
\diagram{\deleted{\begin{tikzcd}[column sep=large,row sep=small]
N_\sigma^*\arrow[rr,"\delta_\sigma"]\arrow[dr,"\varepsilon_\sigma"']&&
\gamma_{\rho,\sigma}\arrow[dl,"\varepsilon_\gamma"]\\
&\{\pm1\}&
\end{tikzcd}}},
\]
\unsure{thus|hence this} commutative\ diagram\ of exact sequences
\[
\tag{43}
\diagram{\begin{tikzcd}[column sep=large]
1\arrow[r]&Z_\sigma\arrow[r]\arrow[d,"\text{epimorphism}"]&N_\sigma^*\arrow[r,"\varepsilon_\sigma"]\arrow[d,"\text{epimorphism}"']&\{\pm1\}\arrow[r]\arrow[d,equal]&1\\
1\arrow[r]&\gamma''_{\rho,\sigma}\arrow[r]&\gamma_{\rho,\sigma}\arrow[r,"\varepsilon_\gamma"]&\{\pm1\}\arrow[r]&1
\end{tikzcd}},\qquad\text{where}
\]
\[
\tag{44}
\gamma''_{\rho,\sigma}=\operatorname{Ker}\bigl(
\diagram{\begin{tikzcd}
\gamma_{\rho,\sigma}\arrow[r,"\varepsilon_\gamma"]&\{\pm1\}
\end{tikzcd}}\bigr).
\]
\mspage{555}

{\small
\begin{center}
555\\
\texttt{https://grothendieck.umontpellier.fr}\\
Université de Montpellier, UMR CNRS 5149, Case courrier 051, Place Eugène Bataillon\\
34095 Montpellier cedex 5 - France
\end{center}
}

\scanpage{557}

\uline{N.B.} The notations $\varepsilon_\gamma,\varepsilon'_\gamma:
\gamma_{\rho,\sigma}\longrightarrow\{\pm1\}$ can lead to confusion,
the first \unsure{referring|pertaining} to the action of a $u$ on $L_\sigma$
(not on $L_\rho$). \unsure{A} slightly heavier notation,
but not leading to confusion, would be
$\varepsilon_{\gamma,\rho},\varepsilon_{\gamma,\sigma}$ --- or even
just $\varepsilon_\rho,\varepsilon_\sigma$, the confusion with the homomorphisms of this name
$\unsure{N'_\rho|N_\rho^*}\longrightarrow\{\pm1\}$,
$N_\sigma\longrightarrow\{\pm1\}$ being harmless when
$\rho^2\ne1,\sigma^2\ne1$, because of the compatibilities (41) (43).

We have thus constructed, with the aid \unsure{of|of the} (1) i.e. of $\mathcal G$ equipped with
$\rho,\sigma,\tau'$, groups \inserted{$\mathcal G^+$,}
$G_{\rho,\sigma},\gamma_{\rho,\sigma},Z_{\rho,\sigma},N'_\rho,N'_\sigma$,
fitting into a commutative diagram of groups
\[
\tag{45}
\diagram{\begin{tikzcd}[column sep=large,row sep=small]
&L_0\simeq SZ_{\rho,\sigma}\arrow[rr]\arrow[ddd,hook]&&Z_{\rho,\sigma}
\arrow[ddd,hook,dashed]\arrow[dddrr,dashed,"\delta_Z"]&&&\\
&&SN_\rho^*\arrow[rr]\arrow[ddl,hook]& &N_\rho^*
\arrow[ddl,hook,dashed]\arrow[ddr,dashed,"\delta_\rho"]&&
\widehat{\mathbb Z}^{*}\times\{\pm1\}\times\{\pm1\}\\
&&SN_\sigma^*\arrow[rr]\arrow[dl,hook,dotted]\arrow[dl]&&N_\sigma^*
\arrow[dl,hook,dotted]\arrow[dl]\arrow[dr,dashed,"\delta_\sigma"]&&\\
1\arrow[r]&\mathcal G^+\arrow[rr,hook,"i_{\mathcal G^+}"]\arrow[dr,hook]&&
G_{\rho,\sigma}\arrow[rr,"\delta_G"]&&\gamma_{\rho,\sigma}\arrow[r]\arrow[uur,"{(\chi_\gamma,\varepsilon_{\gamma,\rho},\varepsilon_{\gamma,\sigma})}"]&1\\
&&\mathcal G\arrow[ur,hook]&&&&
\end{tikzcd}}
\]
where we have set
\[
\tag{46}
\left\{\begin{aligned}
SZ_{\rho,\sigma}&=i_{L_0,Z}(L_0)
=\operatorname{Ker}(\delta_Z:Z_{\rho,\sigma}\longrightarrow\gamma_{\rho,\sigma})\\
SN'_\rho&=\operatorname{Ker}(\delta_\rho:N_\rho^*\longrightarrow\gamma_{\rho,\sigma})
=\operatorname{Cent}_{\mathcal G^+}(\rho)\\
SN'_\sigma&=\operatorname{Ker}(\delta_\sigma:N_\sigma^*\longrightarrow\gamma_{\rho,\sigma})
=\operatorname{Cent}_{\mathcal G^+}(\sigma)
\end{aligned}\right.
\]
\margin{N.B. We can thus also denote them by $SZ_\rho,SZ_\sigma$ instead of the
$SN_\rho^*,SN_\sigma^*$.}
\mspage{556}

{\small
\begin{center}
556\\
\texttt{https://grothendieck.umontpellier.fr}\\
Université de Montpellier, UMR CNRS 5149, Case courrier 051, Place Eugène Bataillon\\
34095 Montpellier cedex 5 - France
\end{center}
}

\scanpage{558}
\mspage{582}

\unsure{in} such a way that we have exact sequences
\[
\tag{47}
\left\{
\diagram{\begin{tikzcd}[column sep=large]
1\arrow[r]&S\unsure{\deleted{Z}N}_\rho^*\arrow[r]&N_\rho^*\arrow[r,"\delta_\rho"]&
\gamma_{\rho,\sigma}\arrow[r]&1\\
1\arrow[r]&SN_\sigma^*\arrow[r]&N_\sigma^*\arrow[r,"\delta_\sigma"]&
\gamma_{\rho,\sigma}\arrow[r]&1
\end{tikzcd}}
\right.
\]
\margin{We have
\[
\tag{47bis}
\left\{\begin{aligned}
\operatorname{Im}\bigl(\unsure{SN_\rho}\longrightarrow\mathcal G^+\bigr)
&=\operatorname{Norm}_{\mathcal G^+}(L_\rho)\\
\operatorname{Im}\bigl(SN_\sigma\longrightarrow\mathcal G^+\bigr)
&=\operatorname{Norm}_{\mathcal G^+}(L_\sigma)
\end{aligned}\right.
\]
(or again, if $\rho^2\ne1$ resp. $\sigma^2\ne1$,
$N_\rho\cap\mathcal G$ and $N_\sigma\cap\mathcal G$)}
mapping to the exact sequence
\[
\diagram{\begin{tikzcd}[column sep=large]
1\arrow[r]&\mathcal G^+\arrow[r]&G_{\rho,\sigma}\arrow[r]&
\gamma_{\rho,\sigma}\arrow[r]&1
\end{tikzcd}},
\]
in a way analogous to (27). Moreover, the homomorphisms
$i_\rho:N_\rho^*\longrightarrow G_{\rho,\sigma}$ and
$i_\sigma:N_\sigma^*\longrightarrow G_{\rho,\sigma}$ have
a marked tendency to be injective (cf. (39)),
which we have marked on diagram (45) by the sign
$\diagram{\begin{tikzcd}[column sep=1.7em]
\phantom{0}\arrow[r,hook,dotted]\arrow[r]&\phantom{0}
\end{tikzcd}}$
(dotted inclusion sign).

We also introduce
\[
\tag{48}
\left\{\begin{aligned}
\chi_\rho&=\chi_\gamma\circ\delta_\rho:
\diagram{\begin{tikzcd}
\unsure{\deleted{Z}N}_\rho^*\arrow[r]&\widehat{\mathbb Z}^{*}
\end{tikzcd}}\\
\chi_\sigma&=\chi_\gamma\circ\unsure{\delta_{\deleted{\rho}\sigma}}:
\diagram{\begin{tikzcd}
\unsure{\deleted{Z}N}_\sigma^*\arrow[r]&\widehat{\mathbb Z}^{*}
\end{tikzcd}}
\end{aligned}\right.,
\]
whose kernels (which contain $S\unsure{\deleted{Z}N}_\rho\deleted{\illegible{1}}$
resp. $SN_\sigma$, and which we \unsure{shall take care not} to confuse with them)
are denoted as is only right by $N_\rho^{*+},N_\sigma^{*+}$.

\uline{N.B.} We should still examine the composites
\[
\tag{49}
\left\{
\diagram{\begin{tikzcd}[column sep=large]
N_\rho^*\arrow[r]&\gamma_{\rho,\sigma}\arrow[r,"\unsure{\varepsilon_{\gamma,\sigma}}"]&\{\pm1\}\\
N_\sigma^*\arrow[r]&\gamma_{\rho,\sigma}\arrow[r,"\varepsilon_{\gamma,\rho}"]&\{\pm1\}
\end{tikzcd}}
\right.,
\]
we shall see that in the ``good cases'', the restrictions of these homomorphisms to
$Z_\rho,Z_\sigma$ respectively are epimorphisms.
\mspage{557}

{\small
\begin{center}
557\\
\texttt{https://grothendieck.umontpellier.fr}\\
Université de Montpellier, UMR CNRS 5149, Case courrier 051, Place Eugène Bataillon\\
34095 Montpellier cedex 5 - France
\end{center}
}